% Expected Attention result narratives — select one after experiment
% Insert into conclusion.tex, replacing the Expected Attention future work bullet

% ============================================================
% SCENARIO A: Expected Attention also produces 60% NIAH
% (Use this if EA does NOT break the NIAH cliff)
% ============================================================

% For conclusion.tex — replace paragraph 2 (eviction causal limitation):
%
% \paragraph{2. The causal limitation is absolute for post-training eviction.}
% Three fundamentally different scoring strategies --- position-based (StreamingLLM),
% attention-based (H\textsubscript{2}O), and statistically-predicted (Expected
% Attention with Gaussian query model) --- all produce identical 60\% NIAH at
% 50\% eviction. This convergence across independent methods establishes that
% eviction-based retrieval degradation is a fundamental information-theoretic
% constraint, not an algorithmic limitation. No post-training scoring function
% operating at prefill time can preserve retrieval for tokens whose importance
% depends on unknown future queries. The only escape route is approaches that
% avoid permanent information loss, such as selective recompute
% (\S\ref{sec:conclusion}).

% For conclusion.tex — replace the Expected Attention future work bullet:
%
% \item \textbf{Selective recompute as the sole path beyond the eviction limit:}
%   We have shown that three independent eviction scoring strategies converge
%   on the same NIAH degradation. Selective recompute --- storing all tokens at
%   ultra-low precision plus token IDs, then recomputing top-$k$ at full precision
%   on the fly --- is the only post-training approach that can preserve retrieval
%   while achieving $>$4$\times$ compression, because it avoids permanent
%   information loss.


% ============================================================
% SCENARIO B: Expected Attention improves NIAH (e.g., 80%+)
% (Use this if EA breaks the NIAH cliff)
% ============================================================

% For conclusion.tex — replace paragraph 2:
%
% \paragraph{2. The causal limitation is partially solvable via statistical modeling.}
% While position-based (StreamingLLM) and attention-based (H\textsubscript{2}O)
% eviction both produce 60\% NIAH at 50\% eviction, Expected Attention
% --- which models future queries as Gaussian-distributed and computes
% closed-form expected attention scores --- achieves [X]\% NIAH at the same
% eviction rate, a [Y]-point improvement. This demonstrates that statistical
% prediction of future attention can partially resolve the causal limitation,
% though the improvement is bounded by the Gaussian assumption: queries that
% deviate significantly from the training distribution (e.g., multi-hop
% reasoning) may still fail.

% For conclusion.tex — replace the Expected Attention future work bullet:
%
% \item \textbf{Expected Attention as NIAH-safe eviction:} Our results show
%   that Gaussian query modeling achieves [X]\% NIAH at [Y]$\times$ compression,
%   partially resolving the causal limitation. Future work should investigate
%   (1) tighter query distribution models (e.g., mixture of Gaussians for
%   multi-intent contexts), (2) per-head adaptive budgets using query
%   self-similarity as an allocation signal, and (3) combining Expected
%   Attention eviction with quantization for multiplicative compression.
